\section{A.L.3: Curvas características}
\sangria{} El objetivo de esta actividad es obtener las diferentes curvas características del transistor NPN utilizando dos fuentes variables que permitan obtener las tensiones y corrientes del capacitor. \\
\sangria{} Se llevó a cabo este circuito:
\begin{center}
\begin{tikzpicture}
	% Paths, nodes and wires:
	\draw (3, 3.25) to[american voltage source, l={$V_1 $}, label distance=0.02cm] (3, 1.5);
	\draw (3.75, 3.5) to[american resistor, l={$R_1$}, label distance=0.02cm] (5.75, 3.5);
	\draw node[npn] (N1) at (6.84, 3.5) {} node[anchor=west] at (N1.text){$BC547C$};
	\draw node[ground] at (3, 1.25) {};
	\draw (3, 0.75) -| (3, 1.5);
	\draw (3, 3.25) -| (3, 3.5) -- (3.75, 3.5);
	\draw (5.75, 3.5) -- (6, 3.5);
	\draw (6.84, 2.73) |- (3, 1.5);
	\draw (6.84, 4.27) to[american resistor, l={$R_2$}, label distance=0.02cm] (6.75, 6.25);
	\draw (9.25, 4.25) to[american voltage source, l={$V_2$}, label distance=0.02cm] (9.25, 2.5);
	\draw (6.75, 6.25) -| (6.75, 6.5) -| (9.25, 4.25);
	\draw (9.25, 2.5) |- (6.84, 1.5);
\end{tikzpicture}
\end{center}
Siendo:
\begin{itemize}
    \item $R_1$ = $R_B$ = $100k\Omega${\small $1/4w$}
    \item $R_2$ = $R_C$ = $560\Omega${\small $1/4w$}
    \item Tanto $V_1$ como $V_2$ son fuentes variables de $~0V$ a $~10V$.
\end{itemize}
\subsubsection{Simulación con LtSpice}
\sangria{} Se implementó el circuito mencionado en el simulador LtSpice para poder observar, haciendo un barrido de tensión del colector $V_{CC}$ y barrido de la tensión de base $V_{BB}$. Con esto, el programa permite graficar la corriente del colector $I_C$ en función de la tensión colector emisor ($V_{CE}$) para las diferentes corrientes de base.

\midTitle{blue}{Notas del documentador:}{configuración LtSpice}
''Al igual que la actividad de laboratorio anterior, se midieron los componentes usados y hemos usado los valores que obtuvimos para llevarlos a la simulación, para visualizar mejor este escenario.''

\subsubsection{Implementación en protoboard}
\begin{center}
    \tikzset{
          hole/.style = {fill = #1, circle, minimum width = 1mm, inner sep = 0mm},
          wire/.style = {draw = #1, line width = 0.2mm},
          label/.style = {rotate = 0, black!50!white, font=\sffamily, anchor = center}
        }
        \begin{tikzpicture}[x=2mm,y=2mm]
          \draw[fill=black!05!white] (-4,-1) rectangle (16,34);
 
          \foreach \yi in {1,...,30}{
            \node[hole=red]  at (-3,\yi){};
            \node[hole=blue] at (-2,\yi){};
            \node[hole=red]  at (14,\yi){};
            \node[hole=blue] at (15,\yi){};

            \foreach \xs in {0,6}{
                \foreach \xi in {1,2,3,4,5}{
                    \node[hole=black] at (\xi+\xs,\yi){};
                }
            }
        }

        \foreach \m/\xi in {A/1,B/2,C/3,D/4,E/5,F/7,G/8,H/9,I/10,J/11}{
            \node[label=0] at (\xi, 0) {\tiny\m};
            \node[label=0] at (\xi,31) {\tiny\m};
        }

        \foreach \n/\yi in {1,5,10,15,20,25,30}{
            \node[label=0] at (-1,\yi) {\tiny\n};
            \node[label=0] at (13,\yi) {\tiny\n};
        }
        \begin{scope}[rotate=180]
        \filldraw[fill=black!60, line width=0.5pt] (-5,-20.75) arc (90:270:0.25cm);
        \filldraw [line width=0.5pt, draw=black, fill=black!60] (-5, -20.75) rectangle (-4.7, -23.25);
        \end{scope}
        \filldraw[fill=brown!60, line width=0.5pt, draw=black] (2.5, 25) rectangle (3.5, 27);
    \filldraw[fill=green, line width=0.25pt, draw=black] (2.5, 25.25) rectangle (3.5, 25.5);
    \filldraw[fill=blue, line width=0.25pt, draw=black] (2.5, 25.6) rectangle (3.5, 25.85);
    \filldraw[fill=brown, line width=0.25pt, draw=black] (2.5, 25.95) rectangle (3.5, 26.2);
    \filldraw[fill=gold, line width=0.25pt, draw=black] (2.5, 26.3) rectangle (3.5, 26.55);
    \draw[line width=1pt, draw=gray!80] (3, 23) -- (3, 25);
    \draw[line width=1pt, draw=gray!80] (3, 27) -- (3, 28);

        \draw[line width=1.5pt, draw=red!90] (4, 28) -- (4, 40) -- (-2.3, 40);
        \filldraw[line width=1.5pt, fill=orange!25] (-1, 40) circle (0.60cm);
        \filldraw[draw=black!90] (-4.25,40) circle (2pt);
        \filldraw[draw=black!90] (2.25,40) circle (2pt);
        \node at (1, 40) {\large\textbf{+}};
        \node at (-3, 39.85) {\large\textbf{-}};
        \node at (-1, 45) {\Large\textbf{$V_{2}$} variable};
        \draw[line width=1.5pt, draw=blue!90] (-2, 30) -- (-2, 35) -- (-4.4, 35)-- (-4.4, 40);
    \filldraw[line width=0.5pt, fill=brown!60, draw=black] (0.5, 21.5) rectangle (-1.5, 22.5);
        \draw[line width=1pt, draw=gray!80] (0.5, 22) -- (2, 22);
        \draw[line width=1pt, draw=gray!80] (-1.5, 22) -- (-3, 22);
    \filldraw[line width=0.25pt, fill=brown, draw=black] (0.25, 21.5) rectangle (0, 22.5);
    \filldraw[line width=0.25pt, fill=black, draw=black] (-0.1, 21.5) rectangle (-0.35, 22.5);
    \filldraw[line width=0.25pt, fill=red, draw=black] (-0.45, 21.5) rectangle (-0.7, 22.5);
    \filldraw[line width=0.25pt, fill=gold, draw=black] (-0.8, 21.5) rectangle (-1.05, 22.5);
    \draw[line width=1pt, draw=blue!60] (1, 21) -- (-1.75, 21);
        \filldraw[line width=1.5pt, fill=orange!25] (-8, 20) circle (0.60cm);
        \filldraw[draw=black!90] (-8,16.75) circle (2pt);
        \filldraw[draw=black!90] (-8,23.25) circle (2pt);
        \node at (-8, 22) {\large\textbf{+}};
        \node at (-8, 17.5) {\large\textbf{-}};
        \node at (-10, 27) {\Large\textbf{$V_{1}$} variable};
        \draw[line width=1.5pt, draw=red] (-7.95, 23.25) -- (-7.95, 24) -- (-3, 24);
        \draw[line width=1.5pt, draw=blue] (-7.95, 16.75) -- (-7.95, 16) -- (-2, 16) ;

        \draw[-latex, very thick, draw=orange, line width=0.75mm] (23, 35) -- (11, 35) -- (9, 22) -- (6.5, 22);
        \node[anchor=east] at (22,37) {\Large $BC547C$};
        \draw[-latex, very thick, draw=blue, line width=0.75mm] (26, 26.5) -- (5, 26.5);
        \node[anchor=east] at (27,28) {\Large $560\Omega$\small $1/4W$};
        \draw[-latex, very thick, draw=violet, line width=0.75mm] (26, 19) -- (2, 19) -- (0, 20.5);
        \node[anchor=east] at (27,20.5) {\Large $100k\Omega$\small $1/4W$};

    \end{tikzpicture}
\end{center}

\subsubsection{Mediciones}
\begin{center}
\begin{tabular}{|
>{\columncolor[HTML]{FFFFC7}}c |cc|cc|ll|}
\hline
\begin{tabular}[c]{@{}c@{}}$V_{CC}$ {[}V{]}\end{tabular} & \multicolumn{2}{c|}{\cellcolor[HTML]{FFCE93}$I_B = 0$}                                & \multicolumn{2}{c|}{\cellcolor[HTML]{FFCE93}$I_B = 10\mu A$}                          & \multicolumn{2}{c|}{\cellcolor[HTML]{FFCE93}$I_B = 15 \mu A$}                                              \\ \hline
0                                                          & \multicolumn{1}{c|}{\cellcolor[HTML]{FFCCC9}$I_C$} & \cellcolor[HTML]{FFCCC9}$V_{CE}$ & \multicolumn{1}{c|}{\cellcolor[HTML]{FFCCC9}$I_C$} & \cellcolor[HTML]{FFCCC9}$V_{CE}$ & \multicolumn{1}{c|}{\cellcolor[HTML]{FFCCC9}$I_C$} & \multicolumn{1}{c|}{\cellcolor[HTML]{FFCCC9}$V_{CE}$} \\ \hline
0,25                                                       & \multicolumn{1}{c|}{}                              &                                  & \multicolumn{1}{c|}{}                              &                                  & \multicolumn{1}{l|}{}                              &                                                       \\ \hline
0,5                                                        & \multicolumn{1}{c|}{}                              &                                  & \multicolumn{1}{c|}{}                              &                                  & \multicolumn{1}{l|}{}                              &                                                       \\ \hline
1                                                          & \multicolumn{1}{c|}{}                              &                                  & \multicolumn{1}{c|}{}                              &                                  & \multicolumn{1}{l|}{}                              &                                                       \\ \hline
2                                                          & \multicolumn{1}{c|}{}                              &                                  & \multicolumn{1}{c|}{}                              &                                  & \multicolumn{1}{l|}{}                              &                                                       \\ \hline
5                                                          & \multicolumn{1}{c|}{}                              &                                  & \multicolumn{1}{c|}{}                              &                                  & \multicolumn{1}{l|}{}                              &                                                       \\ \hline
10                                                         & \multicolumn{1}{c|}{}                              &                                  & \multicolumn{1}{c|}{}                              &                                  & \multicolumn{1}{l|}{}                              &                                                       \\ \hline
\end{tabular}
\end{center}

\begin{center}\begin{tabular}{|
>{\columncolor[HTML]{FFFFC7}}c |cc|cc|}
\hline
\begin{tabular}[c]{@{}c@{}}$V_{CC}$ {[}V{]}\end{tabular} & \multicolumn{2}{c|}{\cellcolor[HTML]{FFCE93}$I_B = 20 \mu A$}                         & \multicolumn{2}{c|}{\cellcolor[HTML]{FFCE93}$I_B = 25\mu A$}                          \\ \hline
0                                                          & \multicolumn{1}{c|}{\cellcolor[HTML]{FFCCC9}$I_C$} & \cellcolor[HTML]{FFCCC9}$V_{CE}$ & \multicolumn{1}{c|}{\cellcolor[HTML]{FFCCC9}$I_C$} & \cellcolor[HTML]{FFCCC9}$V_{CE}$ \\ \hline
0,25                                                       & \multicolumn{1}{c|}{}                              &                                  & \multicolumn{1}{c|}{}                              &                                  \\ \hline
0,5                                                        & \multicolumn{1}{c|}{}                              &                                  & \multicolumn{1}{c|}{}                              &                                  \\ \hline
1                                                          & \multicolumn{1}{c|}{}                              &                                  & \multicolumn{1}{c|}{}                              &                                  \\ \hline
2                                                          & \multicolumn{1}{c|}{}                              &                                  & \multicolumn{1}{c|}{}                              &                                  \\ \hline
5                                                          & \multicolumn{1}{c|}{}                              &                                  & \multicolumn{1}{c|}{}                              &                                  \\ \hline
10                                                         & \multicolumn{1}{c|}{}                              &                                  & \multicolumn{1}{c|}{}                              &                                  \\ \hline
\end{tabular} \end{center}

\subsubsection{Conclusiones}
\midTitle{Green}{Aciertos$\checkmark$}{}
\midTitle{Red}{Dificultades \textbf{\textit{x}}}{}
\midTitle{blue}{Notas del documentador:}{Cómo lo solucionamos}

